\begin{partdivider}{IV}{POUVOIR, VÉRITÉ ET ORGANISATIONS}\end{partdivider}

\bookchap{11}{Qui possède l'intelligence ?}

Nous avons longuement décrit l'IA comme une force de démocratisation, capable de rendre l'expertise accessible au plus grand nombre. Cette vision n'est pas fausse, mais elle est incomplète. Elle occulterait un fait brut, structurant et probablement durable : \textbf{l'abondance d'intelligence à laquelle l'utilisateur final a accès repose sur une rareté extrême des moyens de production de cette intelligence.}

Derrière la fluidité d'un chatbot gratuit ou peu coûteux se cachent des infrastructures colossales : des centres de données de la taille de villes, des milliers de processeurs graphiques (GPU) ultra-spécialisés, des téraoctets de données d'entraînement, et des milliers de chercheurs pointus. Le coût de développement d'un modèle de langage de pointe se chiffre désormais en centaines de millions de dollars, bientôt en milliards.

C'est un fait : la production d'IA à l'état de l'art est l'un des secteurs les plus capitalistiques de l'histoire humaine. Elle concentre les trois facteurs de production les plus rares du moment : le capital financier, le capital technologique (les puces) et le capital cognitif (les talents de recherche).

Cette concentration pose une question politique fondamentale : qui possède l'intelligence artificielle, et qui en contrôle les orientations ?

Le paysage géopolitique actuel dessine une configuration inédite et tendue. Les États-Unis dominent outrageusement la couche applicative et les modèles fondationnels, avec un oligopole de quelques grands acteurs (OpenAI, Google, Meta, Anthropic, Microsoft) qui captent l'essentiel de l'investissement mondial. La Chine, de son côté, poursuit avec une détermination stratégique la maîtrise des modèles et des infrastructures, dans un cadre de souveraineté numérique où l'IA est à la fois un levier de croissance et un instrument de gouvernement. L'Europe, quant à elle, tente de se frayer un chemin par la régulation — avec l'AI Act — et par le soutien à l'écosystème open source, sans disposer pour l'heure des moyens industriels de rivaliser sur les modèles propriétaires de frontière. Mais le vrai point aveugle de cette carte du monde est ailleurs. Il concerne les pays qui ne sont ni producteurs de modèles, ni grands régulateurs, ni même pourvus d'infrastructures de calcul adéquates : une grande partie de l'Afrique, du monde arabe, de l'Asie du Sud-Est ou de l'Amérique latine.

On observe que pour ces pays, l'IA n'est pas une technologie qu'ils produisent ; c'est un service qu'ils importent. Ils sont «consommateurs d'intelligence» dans un système qu'ils ne contrôlent pas. Leurs données ont potentiellement servi à entraîner les modèles du Nord, mais ils n'en maîtrisent ni les usages, ni les biais, ni les orientations. On parle parfois, à juste titre, de risque de «colonialisme de données» ou de «colonialisme algorithmique».

Cela ne signifie pas que l'IA soit nécessairement néfaste pour ces régions. Comme nous l'avons vu, elle peut combler des vides critiques en santé ou en éducation. Mais cela signifie que leur dépendance infrastructurelle s'accroît. Un pays qui sous-traite son intelligence à des serveurs situés à 10~000 kilomètres, régis par le droit d'une autre puissance, perd un peu de son autonomie.

Au sein même des sociétés développées, la question de la propriété se pose avec acuité. L'IA risque-t-elle de réduire ou d'aggraver les inégalités de pouvoir ?

L'hypothèse est la suivante : nous assistons à l'émergence d'une nouvelle classe de pouvoir : les propriétaires de l'infrastructure cognitive. De même que les propriétaires terriens, puis les industriels, puis les pétroliers ont structuré les rapports de force de leurs époques, ceux qui contrôlent les modèles, les données et le \emph{compute} pourraient exercer une influence sans précédent. Car contrôler l'infrastructure de l'intelligence, c'est contrôler une part croissante de la manière dont une société produit, décide, s'informe et se coordonne.

Face à cette concentration, une force de contre-pouvoir émerge cependant : l'open source. Le mouvement qui rend les modèles de langage ouverts, modifiables et téléchargeables par tous (à l'image de Llama par Meta, Mistral par la startup française éponyme, ou des modèles de Hugging Face) est une résistance vitale. Il démocratise la capacité de \emph{fine-tuner} un modèle, de l'adapter à une langue locale, à une culture ou à un usage spécifique sans dépendre d'une API propriétaire.

La tension entre IA propriétaire et IA ouverte est l'une des lignes de faille majeures des prochaines décennies. Si les modèles propriétaires gardent l'avantage de la puissance brute, l'open source permet l'appropriation locale, la transparence et l'innovation par la base. L'issue de cette tension déterminera, pour une large part, le degré de concentration du pouvoir dans la société de l'IA.

Il ne faut cependant pas se leurrer : même avec l'open source, le coût de l'entraînement initial et l'accès aux puces restent des barrières. La véritable question politique n'est pas seulement de savoir si le code est ouvert, mais qui possède les usines qui fabriquent les puces, qui contrôle les câbles sous-marins qui transportent les données, et qui fixe les normes d'interconnexion.

\begin{aiquote}
L'intelligence peut devenir abondante pour l'utilisateur, mais l'infrastructure qui la produit reste profondément hiérarchique. Quelles institutions construirons-nous pour équilibrer ce nouveau rapport de force ?
\end{aiquote}

\bookchap{12}{Une société où chacun peut fabriquer la réalité}

Le chapitre précédent nous a montré que le pouvoir de l'IA réside dans la maîtrise de son infrastructure. Il existe un autre pouvoir, plus diffus, plus insidieux, qui découle directement de la capacité de génération : \textbf{le pouvoir de fabriquer la réalité.}

Jusqu'à une époque récente, produire une image, une voix ou une vidéo nécessitait du matériel, du talent, du temps et — surtout — la mise en œuvre de moyens objectifs. Une photographie laissait une trace physique sur une pellicule. Un enregistrement vocal capturait les vibrations réelles d'une voix. Un journal imprimé portait la mention d'un éditeur responsable. Cette matérialité ne garantissait pas la véracité, mais elle fournissait des points d'ancrage pour la confiance. Nous vivions sous un présupposé hérité de la modernité : \textbf{ce que nous voyons ou entendons, si cela a les apparences de la réalité, a de fortes chances d'être réel.}

L'IA générative brise ce présupposé.

C'est un fait : il est aujourd'hui possible de générer, en quelques secondes et à un coût dérisoire, une image indiscernable d'une photographie, une voix clonée parfaitement, ou une vidéo montrant une personnalité publique prononçant des mots qu'elle n'a jamais dits (les \emph{deepfakes}). Le coût de la contrefaçon de la réalité s'est effondré, tandis que le coût de la vérification, lui, n'a cessé d'augmenter.

L'actualité l'illustre crûment : lors d'élections primaires aux États-Unis début 2024, des appels téléphoniques robotisés utilisant un clone vocal troublant d'un candidat démocrate ont été massivement diffusés pour dissuader les citoyens de voter. La technologie, peu coûteuse et aisément accessible, a prouvé qu'elle pouvait être une arme de déstabilisation électorale à grande échelle, sans qu'aucun acteur humain n'ait eu à prononcer le moindre mot.

Considérons la scène suivante, qui n'a rien de fictif. Quelques semaines avant une élection majeure, une vidéo circule sur les réseaux sociaux. Elle montre un candidat proférant des propos racistes. La vidéo est fausse, générée par une IA. Mais elle est hyper-réaliste. Le camp adverse la dénonce comme un fake, mais cette dénonciation technique n'a pas l'emprise émotionnelle de la vidéo. Le doute s'installe. Le candidat doit s'expliquer, perd du temps, et une partie de l'électorat, sans y croire tout à fait, se dit qu'«il n'y a pas de fumée sans feu».

C'est l'effet bien connu du \emph{firehose of falsehood} (le tuyau d'incendie de mensonges) : l'objectif n'est plus de faire croire un mensonge précis, mais de créer un doute universel. Si tout peut être faux, alors plus rien n'est sûr. La confiance s'effondre.

Notre interprétation est la suivante : l'IA ne crée pas la désinformation — elle est aussi vieille que la politique. Mais elle en change l'économie. Autrefois, désinformer à grande échelle demandait des moyens d'État ou de puissants médias. Aujourd'hui, un petit groupe, voire un individu isolé dans son salon, peut inonder l'espace public de contenus contrefaits adaptés à chaque micro-cible électorale.

Les conséquences institutionnelles sont vertigineuses.

Pour la \textbf{justice}, comment admettre une preuve vidéo si sa manipulation est indécelable à l'œil nu ? Le principe de la preuve par l'image, qui a révolutionné le droit au XX\up{e} siècle, est frappé d'obsolescence technique.

Pour le \textbf{journalisme}, le modèle reposait sur la vérification des sources. Mais comment vérifier une source quand l'IA peut générer un faux site web, de faux profils sociaux, de faux témoignages cohérents entre eux, créant une illusion de consensus complète ?

Pour la \textbf{démocratie}, le délibératif suppose un espace commun de faits avérés sur lesquels débattre d'opinions divergentes. Si les faits eux-mêmes deviennent insaisissables, le débat se disloque. Nous ne débattrons plus de ce qu'il faut faire face à une réalité commune, mais de la nature de cette réalité elle-même.

Face à ce péril, deux types de réponses émergent.

La première est technologique. On tente de créer des outils de détection des contenus générés, ou d'intégrer des signatures invisibles (\emph{watermarking}) et des métadonnées de provenance (comme la norme C2PA) pour certifier l'origine d'une image ou d'un son. C'est indispensable. Mais c'est une course perdue d'avance par la détection : par définition, les modèles génératifs s'améliorent pour corriger les artefacts que les détecteurs identifient. Le \emph{watermarking} ne fonctionne que si les plateformes le respectent et si les utilisateurs prennent la peine de le vérifier — ce qui est peu probable dans le feu de l'actualité.

La seconde réponse est sociale et institutionnelle. Elle consiste à accepter que la preuve technique de la réalité est durablement affaiblie, et à rebâtir la confiance sur d'autres bases. Non plus «je crois ce que je vois», mais «je crois cette source parce que je connais son institution, son historique de fiabilité, et les mécanismes de reddition de comptes auxquels elle est assujettie».

D'où une hypothèse : nous pourrions assister à un retour paradoxal à la confiance humaine. À mesure que la machine sera capable de tout simuler, la garantie de vérité ne viendra plus de la forme du message (l'image, le son) mais de la relation : l'institution qui se porte garante, le journaliste qui a enquêté, le chercheur qui signe de son nom. La valeur de la signature humaine, comme acte d'engagement et de responsabilité, pourrait devenir suprême dans un monde de contenus synthétiques.

Cependant, ce scénario suppose des institutions solides, légitimes et accessibles. Or, c'est précisément la confiance dans les institutions qui est, dans de nombreux pays, au plus bas. Si l'érosion de la confiance technique se conjugue avec une défiance préexistante envers les élites et les médias, alors le risque n'est pas l'adaptation, mais la dislocation : une société où chacun retire dans sa propre version de la réalité, générée à la demande pour confirmer ses préjugés.

\begin{aiquote}
Comment construire une société lorsque nous ne pouvons plus supposer que ce que nous voyons a réellement été produit par un humain ?
\end{aiquote}

C'est là l'un des défis civiques majeurs de notre siècle. Et il ne sera pas résolu par la seule technologie. Il exigera de nouvelles normes sociales, une éducation critique massive (que nous avons évoquée au chapitre 10), et peut-être une redéfinition juridique de la preuve et de l'authenticité.

Ce défi de la vérité n'est pas isolé. Il s'insère dans une transformation plus large de nos organisations. Car l'IA ne modifie pas seulement ce que nous voyons, elle modifie la manière dont nous collaborons, produisons et décidons au sein des entreprises et des administrations. C'est ce que nous devons explorer à présent.

\bookchap{13}{L'entreprise après l'IA}

L'entreprise est, avec l'école, l'institution qui sera le plus profondément remodelée par l'intelligence artificielle. Non pas parce que les entreprises vont «disparaître», mais parce que la logique même de leur organisation, héritée du XX\up{e} siècle, est en contradiction croissante avec les possibilités ouvertes par l'IA.

Pour le comprendre, il faut brièvement rappeler d'où nous venons. L'entreprise moderne, telle qu'elle s'est structurée après la Seconde Guerre mondiale, repose sur un modèle hiérarchique et procédurier. Ce modèle a une justification théorique, théorisée par l'économiste Ronald Coase : l'entreprise existe parce que les coûts de transaction et de coordination sur le marché sont élevés. Il est plus efficace de regrouper des individus dans une hiérarchie, avec des managers pour coordonner leur travail et des processus pour standardiser la production, plutôt que de sous-traiter chaque tâche sur le marché.

Ce modèle a engendré l'organisation que nous connaissons : des pyramides de décision, où l'information remonte difficilement, les instructions redescendent, et des couches de management intermédiaire servent de relais, de filtres et de contrôleurs. Les processus sont conçus comme des chaînes de montage cognitives, où chaque maillon exécute une tâche standardisée.

On observe que l'IA attaque ce modèle sur ses deux fondements : elle réduit drastiquement le coût de coordination, et elle automatise l'exécution des tâches standardisées.

Si un agent IA peut coordonner l'agenda d'une équipe, synthétiser l'avancement de cinq projets, identifier les goulots d'étranglement et proposer des réallocations de ressources en temps réel, une partie de la fonction de management intermédiaire perd sa justification. Si des agents peuvent exécuter de manière autonome la majorité des tâches procédurières (traiter des factures, rédiger des rapports, qualifier des leads, répondre aux questions RH), alors la taille de l'organisation nécessaire pour accomplir un volume de travail donné se contracte.

L'hypothèse est la suivante : l'entreprise de l'ère de l'IA pourrait être plus petite, plus plate et plus réactive. Elle ne fonctionnerait plus sur le schéma \emph{personnes → managers → processus → logiciels}, mais sur le schéma \emph{personnes + agents IA → processus dynamiques}.

Prenons un exemple concret. Aujourd'hui, une startup qui souhaite lancer un produit SaaS (Software as a Service) doit embaucher, au minimum, des développeurs, un designer, un spécialiste marketing, un commercial et un administrateur. C'est la base. Avec l'IA, le scénario change. Un ou deux fondateurs, dotés d'une vision claire, peuvent s'adjoindre des agents de code pour prototyper, des agents de design pour générer les interfaces, des agents de copywriting pour le site web, et des agents d'analyse pour le marché. L'entrepreneur solo, ou la micro-équipe, devient capable de porter un projet d'une complexité autrement inaccessible.

C'est la promesse de la \emph{one-person unicorn} — une entreprise valant un milliard de dollars dirigée par un seul individu. Si ce scénario reste aujourd'hui largement hypothétique, la direction est tracée : \textbf{la productivité par tête explose, ce qui permet de faire plus avec moins de coordination humaine lourde.}

Que deviennent les grandes entreprises, ces dinosaures organisationnels ? Elles ont deux voies devant elles.

La première est la voie de l'optimisation : utiliser l'IA pour faire la même chose, mais plus vite et moins cher. C'est souvent la tentation initiale. On automatise les processus existants, on remplace quelques tâches, on réduit les effectifs. L'organisation, elle, reste hiérarchique. Le risque de cette voie est de ne capter qu'une fraction des gains de l'IA, en préservant des structures dont la lenteur et la rigidité annihilent l'agilité permise par la technologie.

La seconde est la voie de la réinvention : repenser l'organisation \emph{autour} de l'IA, et non pas simplement l'insérer dans l'ancienne. C'est beaucoup plus difficile, car cela suppose de redéfinir les rôles, de supprimer des couches managériales, de faire confiance à des processus pilotés par algorithme, et de reconsidérer ce qui fait la valeur d'un salarié. Dans cette voie, le manager n'est plus un contrôleur qui valide des tâches, mais un architecte de systèmes humains-machine, un définitionneur d'objectifs, un garant d'alignement éthique et stratégique.

Notre interprétation est la suivante : le management de demain ne sera plus l'art de faire faire, mais l'art de faire choisir. Choisir les bons agents, les bonnes données, définir les bons objectifs, et évaluer les résultats dans un environnement où l'exécution est devenue une commodité.

Qu'en est-il du recrutement ? Historiquement, on recrutait pour des compétences techniques : savoir coder, savoir tenir une comptabilité, savoir rédiger. Si ces compétences deviennent des commodités accessibles via IA, que cherche-t-on chez un candidat ? On cherche son jugement, sa capacité à s'adapter, son intelligence relationnelle, sa créativité et son alignement avec la culture de l'entreprise. Le profil du collaborateur idéal se déplace du bon exécutant technique vers le bon décideur situationnel. Cela pourrait, paradoxalement, valoriser des parcours atypiques, des profils libéraux ou humanistes, longtemps marginalisés dans les entreprises techno-centrées.

Il faut cependant nuancer ce scénario. La petite entreprise augmentée est agile, mais elle manque de capital, de réseau et de capacité à absorber les chocs. La grande entreprise a ces actifs, ainsi qu'un accès privilégié aux données massives nécessaires pour entraîner des modèles spécifiques. L'IA pourrait donc, dans certains secteurs, renforcer l'avantage des grands acteurs qui peuvent se payer les meilleurs modèles et intégrer l'IA au cœur de leurs moindres processus, écrasant au passage les PME de taille intermédiaire. Nous assisterions alors non pas à une démocratisation de l'entrepreneuriat, mais à une polarisation entre une myriade de micro-entreprises augmentées (très productives mais vulnérables) et quelques méga-plateformes omniscientes.

L'entreprise après l'IA ne sera donc pas un bloc uniforme. Elle sera le lieu d'une expérimentation organisationnelle intense, dont l'issue dépendra de la capacité des dirigeants à comprendre que l'IA n'est pas un outil de plus, mais une mutation de la manière même dont on coordonne l'action humaine.

Si le travail se transforme, si les organisations se restructurent, si la vérité vacille et si le pouvoir se concentre, alors c'est la société dans son ensemble qui doit se repenser. C'est ce que nous ferons dans notre dernière partie, en explorant les mondes possibles qui s'ouvrent à nous.
